\section{序論}
11月5日07時から6日07時（日本時間）のXでは、セキュリティに特化した高性能モデルと、機械学習の作業をチャットから実行するエージェントが注目を集めた。OpenAIのGPT-6 Cyberは脆弱性研究と安全な利用範囲の両立を掲げ、Hugging FaceのML InternはSlackをモデル評価や学習の操作窓口に変える構想を示した。

研究分野では、モデル出力の出所を証明する暗号技術、シミュレーションと実観測の分布整合、タンパク質配列設計、CAD生成評価など、多様な成果が共有された。オープンウェイトのKimi K3.1 Proや量子化の話題も、モデルを自ら運用するための選択肢への関心を映している。

本号はGrokの07時収集結果12件を記載順に記事化した。投稿の一部は観測窓の直後であり、その境界を明示した。性能値や未確認のモデル公開予測は、原文に記された発表・投稿側の主張として扱う。
